Current AI systems can interpret requests, generate digital outputs, browse the web, and call APIs. Those capabilities are often enough to produce plausible-looking work artifacts. They are not, by themselves, enough to model the full path from ambiguous demand to verified completion. Real work routinely spans clarification, pricing, approval, execution, proof, delivery, and payout. It may also require local access, onsite presence, witnessed handoff, hardware interaction, or field inspection. When those steps are not represented explicitly, the system can still produce polished outputs while failing to capture the actual work needed for completion.

This gap matters more as AI systems move from advisory roles into fulfillment roles. A buyer does not only need a generated answer. The buyer needs a completed outcome, often produced by a mix of humans, agents, tools, and runtimes. Existing stacks tend to fragment this lifecycle across separate surfaces: chat handles intake, a planner generates tasks, a marketplace or directory handles matching, a tracker holds status, and payment systems sit elsewhere. The result is weak continuity. Context drifts across surfaces, approval boundaries blur, and proof of completion is often reduced to narrative summaries rather than durable evidence. Market signals from major work platforms and AI vendors reinforce the importance of this gap: external AI-related work is already monetizing, marketplaces are moving toward fulfillment infrastructure, and enterprise buyers increasingly want systems that connect execution across tools rather than isolated point solutions \cite{upwork_q4_2025,openai_enterprise_ai_2026,microsoft_frontier_suite_2026}.

The problem becomes sharper when fulfillment includes embodied work. Many planners are biased toward the actions they can call directly. If a system can browse a website, produce a checklist, and write a report, it may silently substitute those steps for the real requirement of visiting a site, inspecting an environment, or witnessing a handoff. In other words, the planner overfits to digital executability. This can create outputs that are coherent in language while remaining incomplete in real-world terms.

We argue that a better systems primitive for this setting is a durable request-rooted model. In our approach, one \texttt{Request} persists across intake, routing, \texttt{Commitment}, \texttt{Fulfillment}, proof, and payout. This design treats the request as the canonical work thread rather than as a temporary precursor to separate planner, marketplace, or execution objects. It also allows the system to keep digital and embodied steps inside one accountable lifecycle rather than forcing the physical parts of the work outside the model.

This paper uses Boreal Network as the concrete artifact for exploring that claim. Boreal defines a docs-first and machine-readable contract stack centered on one durable \texttt{Request}, explicit \texttt{Supply}, commercial and approval boundaries through \texttt{Commitment}, execution truth through \texttt{Fulfillment}, proof through \texttt{Artifact}, payment state through \texttt{Transaction}, and append-only history through \texttt{RequestEvent}. Within that model we also define a lead-first orchestration rule: match the lead first, plan the work second, and decompose only when needed.

The paper makes three main contributions. First, it presents a request-rooted object model for mixed human-AI fulfillment that keeps demand, execution, and proof on one durable thread. Second, it identifies a practical failure mode in current agent planning, which we call generative plan collapse: the omission or weakening of required embodied work because the planner is biased toward digitally callable actions. Third, it outlines an artifact-backed evaluation methodology for comparing request-rooted orchestration against chat-first, task-first, and direct-tool baselines on continuity, verification, and false-completion risk.
